\documentclass[11pt,a4paper,fleqn]{article}
\usepackage[T1]{fontenc}
\usepackage[utf8]{inputenc}
\usepackage{lmodern}
\usepackage[czech]{babel}
\usepackage[margin=1.7cm]{geometry}
\usepackage{amsmath,amssymb,mathtools}
\usepackage{xcolor,cancel,enumitem}
\usepackage{tikz,pgfplots}
\pgfplotsset{compat=1.18}
\usepgfplotslibrary{fillbetween}
\usetikzlibrary{arrows.meta,calc,patterns,fit,shapes.geometric,decorations.pathmorphing,decorations.pathreplacing,babel}
\setlength{\mathindent}{1.5em}
\renewcommand{\CancelColor}{\color{red!75!black}}
% značky pro šipky a kroužky (více průchodů pdflatex)
\newcommand{\tm}[2]{\tikz[remember picture, baseline=(#1.base)]\node[inner sep=0pt, outer sep=1pt] (#1) {$\displaystyle #2$};}
\newcommand{\tmt}[2]{\tikz[remember picture, baseline=(#1.base)]\node[inner sep=0pt, outer sep=1pt] (#1) {$\textstyle #2$};}
% zakroužkovaný výraz se jménem: červeně (\krn), modrou propiskou = černě (\krb, tečkovaně \krd)
\newcommand{\krn}[2]{\tikz[remember picture, baseline=(#1.base)]\node[draw=red!75!black, ellipse, inner sep=0.5pt] (#1) {$\displaystyle #2$};}
\newcommand{\krnt}[2]{\tikz[remember picture, baseline=(#1.base)]\node[draw=red!75!black, ellipse, inner sep=0.5pt] (#1) {$\textstyle #2$};}
\newcommand{\krd}[2]{\tikz[remember picture, baseline=(#1.base)]\node[draw, thin, densely dotted, ellipse, inner sep=0pt] (#1) {$\displaystyle #2$};}
\newcommand{\kre}[1]{\tikz[baseline=(k.base)]\node[draw=red!75!black, circle, inner sep=0.3pt, font=\scriptsize] (k) {$#1$};}
\newcommand{\cerv}[1]{{\color{red!75!black}#1}}
\newcommand{\dvoj}[1]{\underline{\underline{#1}}}
% červená závorka pod úsekem (hloubka) a šipka k výsledku
\newcommand{\zav}[3]{\draw[red!75!black, semithick] ([yshift=-2pt]#1.south west) -- ++(0,-#3) -| ([yshift=-2pt]#2.south east);}
\newcommand{\vln}[2]{\tikz[baseline=(o.base)]{\node[inner sep=0pt, outer sep=0pt] (o) {#2};
  \draw[#1, thin, decorate, decoration={snake, amplitude=1pt, segment length=4pt}] ([yshift=-3pt]o.south west) -- ([yshift=-3pt]o.south east);}}
\tikzset{pz/.style={red!75!black, font=\small, inner sep=1pt},
  sipr/.style={-{Stealth}, semithick, red!75!black},
  sipb/.style={-{Stealth}, thin},
  smycka/.style={draw=red!75!black, ellipse, inner sep=-2pt}}
\pagestyle{empty}
\setlength{\parindent}{0pt}
\setlength{\parskip}{4pt}
\newcommand{\vysl}[1]{\boxed{\color{red!75!black}#1}}
\newcommand{\dx}{\,\mathrm{d}x}
\newcommand{\lt}[1]{\lim_{t\to#1}}

\begin{document}
\shorthandoff{-}

{\large Vypočítej obsah (nakresli náčrtek):}

% ---------------------------------------------------------------- a)
\[
\textbf{a)}\quad \tm{a0}{0}\le\tm{ay}{y}\le\tm{af}{\frac{5}{(x-1)^2}}\ ,\hspace{3.2cm} \tm{a2}{2}\le\tm{ax}{x}\le\tm{ai}{\infty}
\begin{tikzpicture}[remember picture, overlay]
\zav{a0}{ay}{8pt}
\zav{ay}{af}{14pt}
\zav{a2}{ax}{8pt}
\zav{ax}{ai}{14pt}
\coordinate (m1) at ($(a0.south)!0.5!(ay.south)+(0,-10pt)$);
\coordinate (m2) at ($(ay.south)!0.5!(af.south)$); \coordinate (m2) at (m2 |- {$(ay.south)+(0,-16pt)$});
\coordinate (m3) at ($(a2.south)!0.5!(ax.south)+(0,-10pt)$);
\coordinate (m4) at ($(ax.south)!0.5!(ai.south)+(0,-16pt)$);
\coordinate (yb) at ($(a0.base)+(0,-1.9cm)$);
\node[anchor=base] (n1) at ($(m1 |- yb)+(-0.7,0)$) {$y\ge0$};
\node[anchor=base] (n2) at ($(m2 |- yb)+(0.55,0)$) {$y\le\dfrac{5}{(x-1)^2}$};
\node[anchor=base] (n3) at ($(m3 |- yb)+(-0.5,0)$) {$x\ge2$};
\node[anchor=base] (n4) at ($(m4 |- yb)+(0.55,0)$) {$x\le\infty$};
\draw[sipb] (m1) -- (n1.north);
\draw[sipb] (m2) -- (n2.north);
\draw[sipb] (m3) -- (n3.north);
\draw[sipb] (m4) -- (n4.north);
\end{tikzpicture}
\]
\vspace{1.3cm}

\hspace*{2.5cm}%
\begin{tikzpicture}
\begin{axis}[axis lines=middle, width=8.5cm, height=6.2cm, clip=false,
  xmin=-2.5, xmax=6.5, ymin=-3, ymax=6, xtick={1,2}, ytick=\empty,
  xticklabel style={font=\small, below right, inner sep=1pt}]
\addplot[pattern=north east lines, draw=none, domain=2:4.5, samples=60] {5/(x-1)^2} \closedcycle;
\addplot[thick, domain=-1.5:0.1, samples=60] {5/(x-1)^2};
\addplot[thick, domain=2:4.8, samples=60] {5/(x-1)^2};
\addplot[thick, domain=1.92:2, samples=20] {5/(x-1)^2};
\draw[dashed] (axis cs:1,-3) -- (axis cs:1,6);
\draw (axis cs:2,-3) -- (axis cs:2,6);
\node[anchor=north west, inner sep=1pt] at (axis cs:2,-3) {$x=2$};
\node[anchor=south west] at (axis cs:3,0.6) {$y=\frac{5}{(x-1)^2}$};
\node[right] at (axis cs:6.5,0) {$y=0$};
\draw[thin] (axis cs:3.6,0.2) -- (axis cs:3.6,-1.6) node[below] {obsah $S=\,?$};
\end{axis}
\end{tikzpicture}

\begin{align*}
S &= \int_2^\infty \krn{h1}{\frac{5}{(x-1)^2}}-\krn{s1}{0}\dx
= \lt{\infty}\int_2^t 5\,(x-1)^{-2}\dx = \lt{\infty}\left[5\,\frac{(x-1)^{-1}}{-1}\right]_2^t =\\[0.4em]
&= \lt{\infty}\left[\frac{-5}{x-1}\right]_2^t = \lt{\infty}\left(\frac{-5}{t-1}-\left(-\frac{5}{2-1}\right)\right)
= \lt{\infty}\Big(\krn{c0}{\frac{-5}{\infty}}+\frac51\Big) = \vysl{5}
\begin{tikzpicture}[remember picture, overlay]
\node[pz, anchor=south] at ([yshift=1pt]h1.north) {Horní $-$ Spodní};
\node[pz, anchor=north] at ([yshift=-1pt]c0.south) {$\overset{\|}{0}$};
\end{tikzpicture}
\end{align*}

\bigskip
% ---------------------------------------------------------------- b)
\[
\textbf{b)}\quad \tm{b0}{0}\le\tm{by}{y}\le\tm{bf}{3^{2x}}\ ,\hspace{3.6cm} \tm{bm}{-\infty}\le\tm{bx}{x}\le\tm{bz}{0}
\begin{tikzpicture}[remember picture, overlay]
\zav{b0}{by}{8pt}
\zav{by}{bf}{14pt}
\zav{bm}{bx}{8pt}
\zav{bx}{bz}{14pt}
\coordinate (m1) at ($(b0.south)!0.5!(by.south)+(0,-10pt)$);
\coordinate (m2) at ($(by.south)!0.5!(bf.south)$); \coordinate (m2) at (m2 |- {$(by.south)+(0,-16pt)$});
\coordinate (m3) at ($(bm.south)!0.5!(bx.south)+(0,-10pt)$);
\coordinate (m4) at ($(bx.south)!0.5!(bz.south)+(0,-16pt)$);
\coordinate (yb) at ($(b0.base)+(0,-1.5cm)$);
\node[anchor=base] (n1) at ($(m1 |- yb)+(-0.85,0)$) {$y\ge0$};
\node[anchor=base] (n2) at ($(m2 |- yb)+(1.2,0)$) {$y\le3^{2x}=\left(3^2\right)^x=9^x$};
\node[anchor=base] (n3) at ($(m3 |- yb)+(-0.7,0)$) {$x\ge-\infty$};
\node[anchor=base] (n4) at ($(m4 |- yb)+(0.55,0)$) {$x\le0$};
\draw[sipb] (m1) -- (n1.north);
\draw[sipb] (m2) -- (n2.north);
\draw[sipb] (m3) -- (n3.north);
\draw[sipb] (m4) -- (n4.north);
\end{tikzpicture}
\]
\vspace{1.2cm}

\hspace*{2.5cm}%
\begin{tikzpicture}
\begin{axis}[axis lines=middle, width=7.5cm, height=5.5cm, clip=false,
  xmin=-3, xmax=1.5, ymin=-1.5, ymax=4.5, xtick={0}, ytick=\empty,
  xticklabel style={font=\small, below right}, axis line style={-}]
\addplot[pattern=north east lines, draw=none, domain=-2.6:0, samples=60] {9^x} \closedcycle;
\addplot[thick, domain=-2.6:0.65, samples=60] {9^x} node[right] {$y=3^{2x}$};
\draw[-{Stealth}, thin] (axis cs:0,0) -- (axis cs:-2.9,0);
\node[below] at (axis cs:-2.9,0) {$-\infty$};
\node[right] at (axis cs:1.5,0) {$y=0$};
\draw[thin] (axis cs:-1.3,0.25) -- (axis cs:-1.6,-0.9) node[below] {obsah $S=\,?$};
\end{axis}
\end{tikzpicture}

\begin{align*}
S &= \int_{-\infty}^0 3^{2x}-0\dx = \lt{-\infty}\int_t^0\krn{c9}{3^{2x}}\dx
= \lt{-\infty}\left[\frac{9^x}{\ln9}\right]_t^0 =\\[0.4em]
&= \lt{-\infty}\frac{9^0}{\ln9}-\frac{9^t}{\ln9} = \frac{1}{\ln9}-\frac{1}{\ln9}\krd{cl}{\lt{-\infty}9^t}
= \frac{1}{\ln9}-\frac{1}{\ln9}\cdot0 = \dvoj{\frac{1}{\ln9}} = \vysl{\frac{1}{2\ln3}}
\begin{tikzpicture}[remember picture, overlay]
\node[pz, anchor=south] at ([yshift=0pt]c9.north) {$\underset{\|}{9^x}$};
\node[anchor=north, inner sep=1pt] (n0) at ([yshift=-1pt]cl.south) {$\overset{\|}{0}$};
\begin{scope}[shift={($(n0.south)+(0,-0.7)$)}, thin]
\draw[-{Stealth}] (0.9,0) -- (-0.9,0) node[below, font=\scriptsize] {$-\infty$};
\draw (0,-0.5) -- (0,0.6);
\draw[thick] (-0.75,0.08) .. controls (-0.3,0.12) and (0,0.18) .. (0.15,0.25) .. controls (0.3,0.33) and (0.45,0.5) .. (0.5,0.6) node[right, font=\small] {$9^t$};
\end{scope}
\end{tikzpicture}
\end{align*}

\newpage
% ---------------------------------------------------------------- c)
\begin{minipage}[t]{0.5\textwidth}
\[
\textbf{c)}\quad
\begin{array}[t]{l}
f\colon\ y=e^x\\[2pt]
g\colon\ y=e^{-x}=\dfrac{1}{e^x}=\left(\krnt{ce}{\frac1e}\right)^x\\[10pt]
\phantom{f\colon\ {}}x=1
\end{array}
\begin{tikzpicture}[remember picture, overlay]
\node[pz, anchor=north west] at ([xshift=-2pt, yshift=-1pt]ce.south east) {$<1$};
\end{tikzpicture}
\]
\end{minipage}%
\begin{minipage}[t]{0.5\textwidth}
\vspace{1pt}
1. nakresli obrázek\\
2. přepiš pomocí nerovnic\\
3. urči obsah ohraničený funkcemi
\end{minipage}

\medskip
\begin{minipage}[t]{0.6\textwidth}
\vspace{0pt}
\textbf{1)}\hspace{0.2cm}%
\begin{tikzpicture}[remember picture, baseline=(current bounding box.center)]
\begin{axis}[axis lines=middle, width=7cm, height=4.9cm, clip=false,
  xmin=-1.8, xmax=2.2, ymin=-1.3, ymax=4, xtick=\empty, ytick=\empty, axis line style={-}]
\addplot[name path=f, thick, domain=-1.6:1.45, samples=60] {exp(x)} node[right] {$y=e^x$};
\addplot[name path=g, thick, domain=-1.4:1.6, samples=60] {exp(-x)};
\node[anchor=south east, inner sep=1pt] at (axis cs:-1.4,4.05) {$y=\left(\frac1e\right)^x=e^{-x}$};
\addplot[pattern=north east lines] fill between[of=f and g, soft clip={domain=0:1}];
\draw (axis cs:1,-1.3) -- (axis cs:1,4) node[above] {$x=1$};
\fill (axis cs:0,1) circle (2pt);
\node[anchor=south west, font=\small] at (axis cs:0,1.05) {$1$};
\node[anchor=north west, font=\small] at (axis cs:1,0) {$1$};
\fill[red!75!black] (axis cs:0,0) circle (1.5pt);
\node[red!75!black, anchor=north east] (z0) at (axis cs:0,-0.05) {$0$};
\draw[thin] (axis cs:0.8,1.2) -- (axis cs:2.2,1.65) node[right] {obsah $S=\,?$};
\end{axis}
\end{tikzpicture}
\end{minipage}%
\begin{minipage}[t]{0.4\textwidth}
\vspace{0.5cm}
\hspace*{0.4cm}Průsečík: $f=g$
\begin{align*}
e^x &= e^{-x} \quad\big/\ln_e\\
\ln_e e^x &= \ln_e e^{-x}\\
x\cdot\krn{l1}{\ln_e e} &= -x\cdot\krn{l2}{\ln_e e}\\[4pt]
x &= -x\\
2x &= 0\\
\tm{x0}{x} &= \tm{x0b}{0}
\begin{tikzpicture}[remember picture, overlay]
\node[pz, anchor=north west] at ([xshift=-2pt]l1.south east) {$=1$};
\node[pz, anchor=north west] at ([xshift=-2pt]l2.south east) {$=1$};
\draw[thin] ([yshift=-2pt]x0.south west) -- ([yshift=-2pt]x0b.south east);
\draw[thin] ([yshift=-3.8pt]x0.south west) -- ([yshift=-3.8pt]x0b.south east);
\draw[sipr, thin] ([yshift=-4pt]x0.south west) to[out=200, in=-60] ([xshift=1pt, yshift=-1pt]z0.south);
\end{tikzpicture}
\end{align*}
\end{minipage}

\medskip
\textbf{2)} Přepis pomocí nerovnic: \
$\begin{array}[t]{c} e^{-x}\le y\le e^{+x}\\[2pt] 0\le x\le1\end{array}$

\begin{align*}
\hspace*{-\mathindent}\textbf{3)}\quad S &= \int_0^1 e^x-e^{-x}\dx = \left[e^x-\frac{e^{-x}}{-1}\right]_0^1
= e^1-\frac{e^{-1}}{-1}-\left(e^0-\frac{e^0}{-1}\right) =\\
&= e+\frac1e-(1+1) = \vysl{e+\frac1e-2}
\end{align*}

\vspace{-2pt}
\tikz\draw[thick, decorate, decoration={snake, amplitude=2.5pt, segment length=22pt}] (0,0) -- (0.98\textwidth,0);
\vspace{-6pt}

% ---------------------------------------------------------------- d)
\begin{minipage}[t]{0.42\textwidth}
\[
\textbf{d)}\quad
\begin{array}[t]{l}
f\colon\ y=\ln x\\
g\colon\ y=\tm{gk}{\ln9}\\
h\colon\ y=\tm{hk}{\ln3}\\
\phantom{h\colon\ {}}x=0
\end{array}
\begin{tikzpicture}[remember picture, overlay]
\node[pz, font=\normalsize, anchor=west] (ko) at ($(gk.east)+(0.6,0.05)$) {konstantní funkce};
\draw[sipr, densely dotted] (ko.west) -- ([xshift=3pt]gk.east);
\draw[sipr, densely dotted] (ko.south west) -- ([xshift=3pt]hk.east);
\end{tikzpicture}
\]
\end{minipage}%
\begin{minipage}[t]{0.58\textwidth}
\vspace{0.9cm}
\begin{tikzpicture}
\begin{axis}[axis lines=middle, width=8cm, height=4.4cm, clip=false,
  xmin=-1, xmax=12, ymin=-2.2, ymax=3.6, xtick=\empty, ytick=\empty]
\addplot[name path=f, thick, domain=0.12:10.3, samples=100] {ln(x)} node[anchor=south west, inner sep=1pt, yshift=6pt, xshift=-4pt] {$f\colon y=\ln x$};
\addplot[name path=g, thick, domain=-0.8:11] {ln(9)} node[right] {$g\colon y=\ln9$};
\addplot[name path=h, thick, domain=-0.8:11] {ln(3)} node[right] {$h\colon y=\ln3$};
\addplot[pattern=north east lines] fill between[of=g and h, soft clip={domain=0:3}];
\addplot[pattern=north west lines, pattern color=red!75!black] fill between[of=g and f, soft clip={domain=3:9}];
\draw[red!75!black, densely dotted] (axis cs:9,0) -- (axis cs:9,{ln(9)});
\draw[red!75!black, densely dotted] (axis cs:3,0) -- (axis cs:3,{ln(3)});
\node[red!75!black, below] at (axis cs:3,0) {$3$};
\node[red!75!black, below] at (axis cs:9,0) {$9$};
\node[below left] at (axis cs:0,0) {$0$};
\fill (axis cs:9,{ln(9)}) circle (1.5pt) node[above left, font=\small] {$P_1$};
\fill (axis cs:3,{ln(3)}) circle (1.5pt) node[below right, font=\small] {$P_2$};
\draw[thin] (axis cs:1.5,1.6) -- (axis cs:1.5,3.3) node[above] {obsah $S_1=\,?$};
\draw[red!75!black, thin] (axis cs:5.5,1.95) -- (axis cs:5.5,2.9) node[above] {obsah $S_2$};
\end{axis}
\end{tikzpicture}
\end{minipage}

\begin{minipage}[b]{0.56\textwidth}
\color{red!75!black}\mbox{\textbf{!}\ Mění se spodní funkce -- PROTO je nutné POČÍTAT}\\
\mbox{}\hfill\mbox{$S_1$, $S_2$ ZVLÁŠŤ\ \textbf{!}}\hspace*{1.2cm}
\end{minipage}\hfill
\begin{tabular}[b]{@{}l@{\ }r@{}l@{}}
Průsečík $P_1$: & $f$ & ${}=g$\\
 & $\ln x$ & ${}=\ln9$\\[5pt]
 & $\tmt{paa}{x}$ & ${}=\tmt{pab}{9}$
\end{tabular}
\quad
\begin{tabular}[b]{@{}l@{\ }r@{}l@{}}
$P_2$: & $h$ & ${}=f$\\
 & $\ln x$ & ${}=\ln3$\\[5pt]
 & $\tmt{pba}{x}$ & ${}=\tmt{pbb}{3}$
\end{tabular}\hspace*{0.3cm}%
\begin{tikzpicture}[remember picture, overlay]
\foreach \a/\b in {paa/pab,pba/pbb} {
  \draw[thin] ([yshift=-2pt]\a.south west) -- ([yshift=-2pt]\b.south east);
  \draw[thin] ([yshift=-3.6pt]\a.south west) -- ([yshift=-3.6pt]\b.south east);
  \coordinate (u) at ([yshift=-4.5pt]\a.south west);
  \node[draw=red!75!black, ellipse, inner sep=1pt, fit=(\a)(\b)(u)] {};
}
\end{tikzpicture}
\vspace{-8pt}
\begin{align*}
\text{Celkový obsah } S &= S_1+S_2 =\\
&= 3\ln3+6-3\ln3 = \vysl{6}
\end{align*}
\vspace{-1.5\baselineskip}
\begin{align*}
S_1 &= \int_0^3 \ln9-\ln3\dx = \Big[\ln9\,x-\ln3\,x\Big]_0^3 = (3\ln9-3\ln3)-(0\ln9-0\ln3) =\\
&= 3(\ln9-\ln3) = 3\ln\frac93 = \dvoj{3\ln3}\\[0.4em]
S_2 &= \int_3^9 \ln9-\krn{lx}{\ln x}^{\,*}\dx = \Big[\ln9\,x-(x\ln x-x)\Big]_3^9 = (9\ln9-9\ln9+9)-(3\ln9-3\ln3+3) =\\
&= 9-3\ln9+3\ln3-3 = 6-3\ln3^2+3\ln3 = 6-6\ln3+3\ln3 = \dvoj{6-3\ln3}
\end{align*}
\hspace*{1.5em}\tikz[remember picture, baseline=(box.base)]\node[draw=red!75!black, thick, rounded corners=2pt, inner sep=4pt] (box) {%
\begin{tabular}{@{}l@{}}
\vln{red!75!black}{Pomocný výpočet:}\\[2pt]
$\displaystyle {}^{*}\int\ln x\dx = \int1\cdot\ln x\dx = \left|\begin{array}{ll} u=\ln x & u'=\frac1x\\ v'=1 & v=x\end{array}\right|
= x\ln x-\int x\cdot\frac1x\dx = \boxed{x\ln x-x}$
\end{tabular}};

\end{document}
